% Stage 8 restored fragment — 2.4.4 Implications for source reconstruction (archive lines 5773-6089)
\subsection{Implications for source reconstruction}
\label{sec:ch244}

A coax-fed pyramidal horn over a ground plane defines the impressed-current and boundary-data tuple technicians declare before Maxwell export is computed \cite{stratton1941}. Maxwell's retarded export to $\Gamma_R$ remains the governing constraint \cite{maxwell1873}.
Exterior radiative uniqueness (Subsection~\ref{sec:ch243}) fixes uniqueness of
\(\mathcal{S}_\omega[\Jimp]\) on the outgoing branch. Subsection~\ref{sec:ch244}
records the inverse-source consequences: field uniqueness does not imply source
identifiability. Uniqueness of \((\E,\H)\) under outgoing conditions does not make
$\mathcal{R}_\omega$ injective---distinct $\Jimp$ can share the same far pattern
\cite{devaney2004}. Source reconstruction is a conditioned inverse, not a direct readout
of coefficients \cite{hansen1988}.
\medskip
\(\mathcal{R}_\omega\) is generally non-injective, and finite observation channels
compress rank further \cite{stratton1941}. Observation null-space
Eq.~\eqref{eq:ch2.observation-null-space}, rank composition
Eq.~\eqref{eq:ch2.rank-composition-law}, and non-radiating relation
Eq.~\eqref{eq:ch2.null-radiation} from Subsection~\ref{sec:ch213} are
presupposed; complete structure of \(\ker\mathcal{R}_\omega\) belongs to
Section~\ref{sec:ch272} \cite{harrington2001}. Measurement classification from
Section~\ref{sec:ch01.5} applies here; angular-momentum observables from
Chapter~\ref{ch:01} are not re-opened \cite{stratton1941}. Field uniqueness versus source identifiability. Forward exterior completion yields

\begin{equation}
\label{eq:ch2.field-source-uniqueness-split}
\underbrace{\mathcal{S}_\omega[\Jimp^{(1)}]=\mathcal{S}_\omega[\Jimp^{(2)}]}_{\text{same impressed data}}
\ \Longrightarrow\
\Jimp^{(1)}=\Jimp^{(2)}
\quad\text{false in general},
\end{equation}
whereas exterior radiative uniqueness (Subsection~\ref{sec:ch243}) gives
\begin{equation}
\label{eq:ch2.field-uniqueness-forward}
\mathcal{S}_\omega[\Jimp^{(1)}]=\mathcal{S}_\omega[\Jimp^{(2)}]
\ \Longrightarrow\
(\E_1,\H_1)=(\E_2,\H_2)
\quad\text{on}\ \mathcal{F}_{\mathrm{SF}}
\end{equation}
for the same boundary class and outgoing branch \cite{stratton1941}. Equations
\eqref{eq:ch2.field-source-uniqueness-split}--\eqref{eq:ch2.field-uniqueness-forward}
separate \emph{state} uniqueness from \emph{excitation} uniqueness: many admissible
currents can produce the same radiative exterior field once reactive degrees of
freedom are quotiented \cite{harrington2001}. Inverse problems therefore target
equivalence classes of \(\Jimp\), not individual source distributions
\cite{stratton1941}.
The radiative export map from Subsection~\ref{sec:ch201} is
\begin{equation}
\label{eq:ch2.radiation-forward-map}
\mathcal{R}_\omega:\mathcal{J}_\Omega\longrightarrow\mathcal{F}_{\mathrm{rad}},
\qquad
\mathcal{R}_\omega[\Jimp]=\Pi_{\mathrm{rad}}\,\mathcal{S}_\omega[\Jimp],
\end{equation}
restating Eq.~\eqref{eq:ch2.radiation-operator-def} after outgoing radiation condition
Eqs.~\eqref{eq:ch2.outgoing-radiation-refinement}--\eqref{eq:ch2.outgoing-resolvent-unique}
\cite{stratton1941}. The source-side null space recorded in
Eq.~\eqref{eq:ch2.null-radiation} is
\begin{equation}
\label{eq:ch2.ker-R-omega}
\ker\mathcal{R}_\omega
=
\left\{
\Jimp\in\mathcal{J}_\Omega:
\Pi_{\mathrm{rad}}\,\mathcal{S}_\omega[\Jimp]=\mathbf{0}
\right\},
\end{equation}
generally infinite-dimensional on open domains \cite{harrington2001}.
Section~\ref{sec:ch272} develops the complete null-space analysis; Subsection~\ref{sec:ch244}
records the reconstruction consequences of Eq.~\eqref{eq:ch2.ker-R-omega}
without constructing non-radiating bases \cite{stratton1941}. Elements of
\(\ker\mathcal{R}_\omega\) are invisible to far-zone radiative observables yet may
carry substantial near-zone reactive energy through \(\mathcal{S}_\omega\)
\cite{harrington2001}.
\begin{figure}[t]
\centering
\ChOneFigBlock{%
\begin{tikzpicture}[ch1fig, node distance=5mm and 7mm]
 \node[ch1box] (j) {$\mathcal{J}_\Omega$};
 \node[ch1box, right=7mm of j] (s) {$\mathcal{S}_\omega$};
 \node[ch1box, right=7mm of s] (f) {$\mathcal{F}_{\mathrm{SF}}$};
 \node[ch1box, right=7mm of f] (r) {$\mathcal{F}_{\mathrm{rad}}$};
 \node[ch1box, below=5mm of j] (ker) {$\ker\mathcal{R}_\omega$};
 \draw[ch1flow] (j) -- (s) -- (f) -- node[above,font=\scriptsize] {$\Pi_{\mathrm{rad}}$} (r);
 \draw[ch1flow, dashed] (ker.south) -- (j.north);
\end{tikzpicture}%
}
\caption{Forward chain for reconstruction: \(\mathcal{S}_\omega\) is single-valued on
\(\mathcal{F}_{\mathrm{SF}}\) after exterior radiative uniqueness (Subsection~\ref{sec:ch243});
\(\mathcal{R}_\omega\) collapses many sources to the same radiative output.}
\label{fig:ch2.reconstruction-forward-chain}
\end{figure}
Observation composition and inverse problem. Laboratory readouts compose with $\mathcal{R}_\omega$ before any source inversion
is attempted: pattern probes and near-field scanners observe a finite-rank image of impressed
currents, not $\Jimp$ itself \cite{hansen1988}. From
Eq.~\eqref{eq:ch2.radiation-operator-composition} and finite-rank observation
Eq.~\eqref{eq:ch2.finite-rank-observation},
\begin{equation}
\label{eq:ch2.reconstruction-observation-map}
\mathcal{M}_\omega
=
\mathcal{O}_{\mathrm{meas}}\circ\mathcal{O}_{\mathrm{far}}\circ\mathcal{R}_\omega
=
\mathcal{O}_{\mathrm{meas}}\circ\Lambda_{\mathrm{rad}}:
\mathcal{J}_\Omega\longrightarrow\mathbb{C}^{N},
\end{equation}
with \(\Lambda_{\mathrm{rad}}=\mathcal{O}_{\mathrm{far}}\circ\mathcal{R}_\omega\)
\cite{stratton1941}. Given measurement vector
\(\mathbf{y}\in\mathbb{C}^{N}\), the inverse source problem seeks
\(\Jimp\in\mathcal{J}_\Omega\) such that
\begin{equation}
\label{eq:ch2.inverse-source-equation}
\mathcal{M}_\omega[\Jimp]=\mathbf{y},
\qquad
\mathcal{D}_{J}[\Jimp]=1,
\end{equation}
subject to admissibility Eq.~\eqref{eq:ch2.source-admissibility}
\cite{harrington2001}. Equation~\eqref{eq:ch2.inverse-source-equation} is ill-posed
whenever \(\ker\mathcal{M}_\omega\neq\{\mathbf{0}\}\), which occurs if either
\(\ker\mathcal{R}_\omega\neq\{\mathbf{0}\}\) or \(\ker\mathcal{O}_{\mathrm{meas}}\)
is nontrivial on \(\operatorname{ran}\mathcal{R}_\omega\)
\cite{stratton1941}. Rank cap Eq.~\eqref{eq:ch2.rank-composition-law} applies directly:
\begin{equation}
\label{eq:ch2.reconstruction-rank-cap}
\mathrm{rank}(\mathcal{M}_\omega)
\le
\min\big\{N,\,\dim\mathcal{R}(\mathcal{R}_\omega)\big\},
\end{equation}
so no measurement protocol with \(N\) channels resolves more than \(N\) independent
radiative degrees of freedom \cite{harrington2001}. The composed null space satisfies
\begin{equation}
\label{eq:ch2.reconstruction-null-space}
\ker\mathcal{M}_\omega
\supseteq
\ker\mathcal{R}_\omega
\ \cup\
\mathcal{R}_\omega^{-1}\!\big(\ker\mathcal{O}_{\mathrm{meas}}\big),
\end{equation}
with set inclusion on admissible sources \cite{stratton1941}. Identifiable sources
modulo null space form the quotient
\begin{equation}
\label{eq:ch2.source-identifiability-quotient}
\mathcal{J}_\Omega\big/\ker\mathcal{M}_\omega
\ \cong\
\operatorname{ran}\mathcal{M}_\omega
\subset\mathbb{C}^{N},
\end{equation}
so reconstruction recovers at best a finite-dimensional image, not a unique
\(\Jimp(\mathbf{r})\) \cite{harrington2001}. Equivalent-current models that expand
\(\Jimp\) on auxiliary surfaces inherit the same non-uniqueness: distinct surface
distributions related by \(\ker\mathcal{R}_\omega\) produce identical far-zone
patterns \cite{stratton1941}. Boundary-element inversion therefore
requires regularization or additional constraints beyond Maxwell consistency alone
\cite{harrington2001}.
\begin{figure}[t]
\centering
\ChOneFigBlock{%
\begin{tikzpicture}[ch1fig, node distance=5mm and 7mm]
 \node[ch1box] (j1) {$\Jimp^{(1)}$};
 \node[ch1box, right=7mm of j1] (j2) {$\Jimp^{(2)}$};
 \node[ch1box, right=7mm of j2] (plus) {$\Jimp^{(1)}+\delta\Jimp$};
 \node[ch1box, below=5mm of plus] (y) {Same $\mathbf{y}$};
 \node[ch1box, below=5mm of j1] (delta) {$\delta\Jimp\in\ker\mathcal{M}_\omega$};
 \draw[ch1flow] (j1.east) -- (j1.east |- plus.west);
 \draw[ch1flow] (j1.east |- plus.west) -- (plus.west);
 \draw[ch1flow] (delta.east) -- (delta.east |- plus.west);
 \draw[ch1flow] (delta.east |- plus.west) -- (plus.west);
 \draw[ch1flow] (plus.south) -- (plus.south |- y.north);
 \draw[ch1flow] (plus.south |- y.north) -- (y.north);
\end{tikzpicture}%
}
\caption{Non-uniqueness: adding \(\delta\Jimp\in\ker\mathcal{M}_\omega\) leaves
measurement vector \(\mathbf{y}\) unchanged.}
\label{fig:ch2.reconstruction-null-class}
\end{figure}
~\eqref{eq:ch2.observation-null-space} pulled back to source space through
Eq.~\eqref{eq:ch2.reconstruction-observation-map} \cite{stratton1941}. When Eq.~\eqref{eq:ch2.inverse-source-equation} is ill-posed, stable inversion uses
a penalized least-squares or Tikhonov template on admissible sources:
\begin{equation}
\label{eq:ch2.tikhonov-source-reconstruction}
{\Jimp}_{\mathrm{reg}}
=
\arg\min_{\Jimp\in\mathcal{J}_\Omega}
\left\{
\|\mathcal{M}_\omega[\Jimp]-\mathbf{y}\|^{2}
+\alpha\,\mathcal{P}[\Jimp]
\right\},
\end{equation}
with regularization parameter \(\alpha>0\) and penalty
\(\mathcal{P}\) enforcing smoothness, minimum norm, or sparsity on an equivalent
surface \cite{harrington2001}. Equation
\eqref{eq:ch2.tikhonov-source-reconstruction} does not remove
\(\ker\mathcal{M}_\omega\); it selects one representative from the affine coset
\(\mathcal{M}_\omega^{-1}(\mathbf{y})\) \cite{stratton1941}. The choice of
\(\mathcal{P}\) is physics-adjacent modeling, not a Maxwell identity; different
penalties yield different \({\Jimp}_{\mathrm{reg}}\) with identical radiative
predictions when \(\delta\Jimp\in\ker\mathcal{R}_\omega\) \cite{harrington2001}. Transport and angular reports attach to \(\mathcal{R}_\omega[\Jimp]\), not to raw
\(\mathcal{S}_\omega[\Jimp]\), per Eq.~\eqref{eq:ch2.transport-qualified-radiation}
\cite{stratton1941}. regime qualification Eq.~\eqref{eq:ch2.regime-operator-screen} criteria
whether \(\mathcal{R}_\omega[\Jimp]\) lies in the propagating sector on which
far-zone observables were designed:
\begin{equation}
\label{eq:ch2.reconstruction-regime-criterion}
\mathcal{D}_{J,\mathrm{reg}}=1
\ \Longrightarrow\
\text{inverse claims target}\ \mathcal{R}_\omega[\Jimp],
\ \text{not reactive near fields},
\end{equation}
with \(\mathcal{D}_{J,\mathrm{reg}}\) from Eq.~\eqref{eq:ch1.regime-qualified-decision}
\cite{stratton1941}. Attempting to reconstruct \(\Jimp\) from
near-zone reactive samples without regime qualification conflates storage with export,
the laboratory error flagged in Section~\ref{sec:ch01.2.4} \cite{stratton1941}. Outgoing
uniqueness from exterior radiative uniqueness (Subsection~\ref{sec:ch243}) stabilizes the
\emph{field} side of Eq.~\eqref{eq:ch2.inverse-source-equation}; it does not promote
evanescent-dominated states to identifiable radiative sources
\cite{harrington2001}.
Formal adjoint back-projection provides a reconstruction \emph{template} without
claiming uniqueness:
\begin{equation}
\label{eq:ch2.adjoint-backprojection}
\widetilde{\Jimp}
=
\mathcal{M}_\omega^{\dagger}\,\mathbf{y},
\end{equation}
where \(\mathcal{M}_\omega^{\dagger}\) is the Hilbert-space adjoint on declared
domains \cite{stratton1941}. Equation
\eqref{eq:ch2.adjoint-backprojection} is useful for sensitivity analysis and
iterative inversion but generally satisfies
\(\mathcal{M}_\omega[\widetilde{\Jimp}]\neq\mathbf{y}\) when
\(\ker\mathcal{M}_\omega\neq\{\mathbf{0}\}\) \cite{harrington2001}. Consistent
reconstruction requires projecting onto \(\operatorname{ran}\mathcal{M}_\omega^{\dagger}\)
or applying Eq.~\eqref{eq:ch2.tikhonov-source-reconstruction}
\cite{stratton1941}. Full integral representations for
Eq.~\eqref{eq:ch2.adjoint-backprojection} belong to Section~\ref{sec:ch251}; they are
not repeated here \cite{harrington2001}.
Equivalent sources and Huygens surfaces. Equivalent-current inversion on a closed surface \(\Gamma_{\mathrm{eq}}\) seeks
\(\mathbf{J}_{\mathrm{eq}}\) on \(\Gamma_{\mathrm{eq}}\) such that exterior fields
match measured data \cite{stratton1941}. Uniqueness of exterior
\(\mathcal{S}_\omega\) from exterior radiative uniqueness (Subsection~\ref{sec:ch243})
fixes the \emph{field} on \(\mathcal{F}_{\mathrm{SF}}\) but not
\(\mathbf{J}_{\mathrm{eq}}\): any non-radiating tangential distribution on
\(\Gamma_{\mathrm{eq}}\) can be added without changing far-zone samples
\cite{harrington2001}. The operator form is
\begin{equation}
\label{eq:ch2.equivalent-source-ambiguity}
\mathbf{J}_{\mathrm{eq},1}-\mathbf{J}_{\mathrm{eq},2}
\in
\ker\mathcal{R}_\omega\big|_{\Gamma_{\mathrm{eq}}},
\end{equation}
with restriction to surface currents in the trace class of
Subsection~\ref{sec:ch211} \cite{stratton1941}. Numerical equivalence in
Eq.~\eqref{eq:ch2.discrete-green-equivalence} inherits the same null-space structure
at finite matrix rank \cite{harrington2001}.
Near-field versus far-field data. Near-field scanning samples \(\mathcal{S}_\omega[\Jimp]\) on a surface enclosing
\(\Omega_s\), not \(\mathcal{R}_\omega[\Jimp]\) alone \cite{stratton1941}. Because
\(\mathcal{S}_\omega\) is injective on outgoing fields but many sources share the
same \(\mathcal{R}_\omega\) image, near-field data can \emph{reduce} null-space
dimension relative to far-zone-only inversion without eliminating
\(\ker\mathcal{R}_\omega\) in general \cite{harrington2001}. Outgoing radiation condition still
required: two incoming/outgoing superpositions with identical near traces on a finite
surface need not agree in the far zone unless exterior radiative uniqueness (Subsection~\ref{sec:ch243})
applies to the full exterior problem \cite{stratton1941}. Near-field
inversion therefore couples Subsection~\ref{sec:ch243} field uniqueness with
Subsection~\ref{sec:ch244} source non-uniqueness \cite{harrington2001}. Equations~\eqref{eq:ch2.inverse-source-equation}--\eqref{eq:ch2.tikhonov-source-reconstruction}
require linear Maxwell theory, fixed \(\omega\), admissible
Eq.~\eqref{eq:ch2.source-admissibility}, and outgoing forward completion
Eq.~\eqref{eq:ch2.outgoing-resolvent-unique} \cite{stratton1941}.
Nonlinear, time-varying, or dispersive inversion breaks the harmonic reduction of
Subsection~\ref{sec:ch223}. Single-frequency inversion cannot resolve broadband
excitation without superposition across \(\omega\) \cite{stratton1941}. Loss improves
conditioning of Eq.~\eqref{eq:ch2.inverse-source-equation} by damping resonances but
does not make \(\mathcal{R}_\omega\) injective on all of \(\mathcal{J}_\Omega\)
\cite{harrington2001}. Complete null-space bases and NR source examples belong to
Section~\ref{sec:ch27} \cite{stratton1941}.
As a worked example, with \(N\) probe channels, Eq.~\eqref{eq:ch2.reconstruction-rank-cap} limits
recoverable modal content to rank \(N\) regardless of mesh density
\cite{stratton1941}. Refining a boundary-element mesh without adding independent
measurements does not resolve components in
\(\ker\mathcal{O}_{\mathrm{meas}}\circ\Lambda_{\mathrm{rad}}\)
\cite{harrington2001}. This is the operator reading of
Eq.~\eqref{eq:ch2.observation-null-space}, distinct from discretization error alone
\cite{stratton1941}. As a worked example, a source with \(\mathcal{D}_{J,\mathrm{reg}}=0\) may produce strong near-zone
\(\mathcal{S}_\omega[\Jimp]\) while \(\mathcal{R}_\omega[\Jimp]\approx\mathbf{0}\)
\cite{stratton1941}. Inverting near samples without
Eq.~\eqref{eq:ch2.reconstruction-regime-criterion} assigns spurious radiative labels to
storage-dominated fields \cite{harrington2001}. regime qualification
Eq.~\eqref{eq:ch2.regime-operator-screen} must precede angular-transport claims on
reconstructed sources \cite{stratton1941}. \medskip
\noindent
Section~\ref{sec:ch24} closes with the excitation-to-flux sequence fixed for open regions:
Sommerfeld asymptotics and outgoing class
Eqs.~\eqref{eq:ch2.sommerfeld-outgoing-class}--\eqref{eq:ch2.sommerfeld-field-condition}
(Subsection~\ref{sec:ch241}), outgoing selector and resolvent branch
Eqs.~\eqref{eq:ch2.outgoing-selector-def}--\eqref{eq:ch2.outgoing-solution-operator}
(Subsection~\ref{sec:ch242}), exterior radiative uniqueness
Eq.~\eqref{eq:ch2.exterior-uniqueness-theorem} and resolvent injectivity
Eq.~\eqref{eq:ch2.exterior-resolvent-injective} (Subsection~\ref{sec:ch243}), and
inverse reconstruction limits
Eqs.~\eqref{eq:ch2.inverse-source-equation}--\eqref{eq:ch2.ker-R-omega}
(Subsection~\ref{sec:ch244}) \cite{stratton1941}. Green
kernels from Section~\ref{sec:ch23} are now tied to a single physical branch;
\(\mathcal{S}_\omega\) is single-valued on \(\mathcal{F}_{\mathrm{SF}}\), while
\(\mathcal{R}_\omega\) and composed measurement readouts remain rank-deficient on
\(\mathcal{J}_\Omega\) \cite{harrington2001}. Section~\ref{sec:ch25} evaluates
radiation integrals, flux meters, and near/intermediate/far partitions on the
qualified propagating subspace; Section~\ref{sec:ch27} completes non-radiating
null-space structure \cite{stratton1941}. Regime qualification
Eq.~\eqref{eq:ch2.regime-operator-screen} and transport criterion
Eq.~\eqref{eq:ch2.transport-qualified-radiation} remain mandatory before
angular-transport claims on operator outputs \cite{stratton1941}.
Outgoing-wave selection fixes the exterior resolvent branch assumed throughout radiation
integrals; realizability limits---which accessible modes finite geometry can
excite---are established in \ref{ch:04} (\emph{Source $\mathcal{S}_{\mathrm{real}}$ and Spectral Truncation}).
